\section{问题分析}

本题以镇龙乡标准测试场景为背景，围绕“运输组批—多架次调度—通信中继协同—任务分区”四个递进环节展开。四个子问题在物理规则（附录 2 时间与能耗）、通信规则（附录 3 链路判定）、对象编号与计算口径上完全统一，并逐层继承数据与中间结果。总体技术路线如图~\ref{fig:roadmap} 所示。

\begin{figure}[htbp]
  \centering
  \includegraphics[width=0.86\textwidth]{roadmap.png}
  \caption{全文技术路线图}
  \label{fig:roadmap}
\end{figure}

\subsection{问题一的分析}

问题一在不考虑实体无人机与共享电池调度的前提下，研究单服务区直接往返（$O01\to S_i\to O01$）的运输能力与货箱组批。其核心是两个耦合的子任务：一是确定各机型在各服务区的\textbf{最大安全载荷}，二是在载荷、体积与返航安全能量余量约束下完成\textbf{货箱组批}并优化。

最大安全载荷问题本质是由返航安全能量余量约束隐式定义的临界载荷求解：随载荷增大，等效航程缩短、单架次能耗上升，当能耗恰好触及安全余量边界时得到能量受限载荷，与机型最大载货质量取较小值即为最大安全载荷。由于能耗对载荷单调，可用一维求根（brentq）稳定求解。组批问题是典型的\textbf{装箱/集合划分}问题：货箱不可拆分、每箱仅安排一次，需同时满足质量与体积两类背包约束，并以架次数、总能耗、累计作业时间为多目标。本文先枚举可行批次，再比较三种方法：线性加权集合划分 MILP、理想点 Chebyshev 集合划分 MILP，以及分服务区 NSGA-II 与 Minkowski 前沿合并。前两种给出精确集合划分解，第三种提供进化搜索得到的非支配方案，用以检验不同求解路径下的代表解与权衡关系。

\subsection{问题二的分析}

问题二放开“单点往返”限制，每个架次可多点访问，并引入实体机队（A/B/C 三型 8 架）、共享电池数量、两阶段充电周转以及医疗物资期望送达、首批保障截止等时限约束，是一个\textbf{带资源约束与时间窗的多行程车辆路径问题（MTVRP）}。本文比较全局分组编码 NSGA-II、路线构造后触发 NSGA-II 修复的混合算法，以及贪心合并基线；三者分别代表全局搜索、可行性修复和低成本构造。代码包中的 p2-2 复用 q2-1 的 NSGA-II 核心，仅以较小种群和代数复跑，属于参数复跑而非独立算法，故不列入主比较表。评价目标统一为“加权迟到量—完工时间—总能耗—架次数”。

\subsection{问题三的分析}

问题三在问题二基础上叠加\textbf{连续通信约束}：运输无人机在爬升、巡航、下降及投送全程必须保持与固定网关 $G01$ 的直连，或由一架中继无人机提供“运输—中继—$G01$”两段链路。通信状态依赖自由空间损耗、双向链路预算与 $30\,$m DEM 视线遮挡。本文将全局分组编码 NSGA-II 定为高复杂度主解：染色体搜索货箱分组，解码阶段据此构造运输架次，并通过联合离散事件仿真评估运输与中继资源调度；中继悬停点由通信缺口与候选点规则确定，并非独立搜索基因。贪心运输路线加中继覆盖分配作为基础解。代码包中的 \texttt{p3-3pro} 固定上游运输路线后优化多站访问顺序与中继点，\texttt{p3-4pro} 进一步考察链路损耗扰动，二者搜索范围和输入链不同，作为增强性参考而不与主表混合排名。

\subsection{问题四的分析}

问题四固定问题三基础解的路线、通信关系和 must-link 约束，研究\textbf{任务分区}对资源配置的影响。本文用完全枚举作为高复杂度精确主解，用局部搜索作为基础解，二者在同一 $13$ 个原子任务块和同一资源核算器上比较；通信风险感知的 `p4-3pro/p4-4pro` 属于另一条 `p3-3pro/p3-4pro` 路线链，单独作为扩展方向，不并入本题主结论。
